\chapter{The screen: coordinates, colour, shapes}

\section{Where is (0,0)?}

Top left corner. X grows to the right, Y grows \emph{downwards}.

\begin{center}
\begin{tikzpicture}[scale=0.95,every node/.style={font=\sffamily\small}]
  \fill[rlight] (0,0) rectangle (9,-4.6);
  \draw[codeframe] (0,0) rectangle (9,-4.6);
  \foreach \x in {1,2,...,8} \draw[codeframe] (\x,0) -- (\x,-4.6);
  \foreach \y in {1,2,3,4} \draw[codeframe] (0,-\y) -- (9,-\y);

  \fill[rink] (0,0) circle (2.2pt);
  \node[above right,font=\sffamily\footnotesize] at (0.1,0.05) {(0,0)};

  \draw[->,rred,line width=1.3pt] (0,0) -- (3,0) node[above,pos=0.8] {X grows right};
  \draw[->,rgreen,line width=1.3pt] (0,0) -- (0,-2.6) node[right,pos=0.85,align=left] {Y grows\\ \textbf{down}};

  \fill[rblue] (5,-2) rectangle (7.4,-3.4);
  \fill[rink] (5,-2) circle (2pt);
  \node[above,font=\sffamily\footnotesize,text=rink] at (5,-1.95) {(x, y)};
  \draw[<->,rgrey] (5,-3.7) -- node[below,font=\sffamily\footnotesize] {width} (7.4,-3.7);
  \draw[<->,rgrey] (7.7,-2) -- node[right,font=\sffamily\footnotesize] {height} (7.7,-3.4);

  \node[below right,font=\sffamily\footnotesize,text=rgrey] at (0.1,-4.75) {DrawRectangle(x, y, width, height, color)};
\end{tikzpicture}
\end{center}

If you did maths at school, Y pointing down feels wrong. It comes from how
screens were scanned, line by line from the top, and every 2D graphics API on
earth inherited it.

\gotcha{In 3D (chapter 9) Y goes back to pointing \emph{up}. Same letter,
opposite direction, in the same program, sometimes on the same screen. When
something moves the wrong way vertically, this is the first thing to check:
which of the two worlds am I in?}

Two helpers give you the current window size, and you should prefer them over
hardcoding the numbers you passed to \lstinline|InitWindow|, because the window
can be resized:

\begin{lstlisting}
int w = GetScreenWidth();
int h = GetScreenHeight();
DrawText("centred-ish", w/2 - 60, h/2, 20, DARKGRAY);
\end{lstlisting}

\section{Corners and centres}

\lstinline|DrawRectangle| takes the \textbf{top-left corner}.
\lstinline|DrawCircle| takes the \textbf{centre}. That inconsistency has burned
every raylib beginner, so hold on to it:

\begin{lstlisting}
DrawRectangle(100, 100, 80, 40, MAROON);   // corner at (100,100), 80x40
DrawCircle(100, 100, 30, DARKBLUE);        // CENTRE at (100,100), radius 30
\end{lstlisting}

If you want a rectangle centred at a point, you subtract half its size yourself:

\begin{lstlisting}
DrawRectangle(cx - w/2, cy - h/2, w, h, MAROON);
\end{lstlisting}

You will write that line a thousand times. In the 3D chapters it comes back as
\lstinline|BoxFromCenter()|, which is the same idea with one more axis.

\section{Rectangle, the struct you will live in}

\begin{lstlisting}
typedef struct Rectangle {
    float x;        // LEFT edge
    float y;        // TOP edge
    float width;
    float height;
} Rectangle;
\end{lstlisting}

Many functions take one of these instead of four loose numbers, and the
collision functions all do. Get used to building them inline:

\begin{lstlisting}
Rectangle player = { 100, 100, 40, 40 };
DrawRectangleRec(player, DARKBLUE);
DrawRectangleLinesEx(player, 2, BLACK);
\end{lstlisting}

\section{Colour}

A \lstinline|Color| is four bytes, each 0 to 255:

\begin{lstlisting}
typedef struct Color { unsigned char r, g, b, a; } Color;

Color myOrange = { 230, 120, 20, 255 };          // opaque
Color ghost    = { 230, 120, 20, 90 };           // mostly transparent
\end{lstlisting}

\lstinline|a| is alpha: 255 is solid, 0 is invisible. raylib ships about
twenty-five named colours --- \lstinline|RAYWHITE|, \lstinline|DARKGRAY|,
\lstinline|MAROON|, \lstinline|SKYBLUE|, \lstinline|GOLD|, \lstinline|BLANK|
(fully transparent) and friends --- and one very handy helper:

\begin{lstlisting}
DrawRectangle(10, 10, 100, 100, Fade(BLUE, 0.5f));   // BLUE at 50% alpha
\end{lstlisting}

\lstinline|Fade| returns a copy of the colour with the alpha scaled. Use it for
shadows, overlays, ghosts and dimmed HUDs instead of writing the struct out.

\gotcha{When you write a colour inline you need the cast, because C needs to
know what the braces mean: \lstinline|DrawRectangle(0,0,10,10,(Color){20,20,25,255})|.
Forgetting the \lstinline|(Color)| gives you a confusing compiler error about
initialisation.}

\section{The order of your draw calls is the layer order}

There is no z-index, no layer system and no sorting. Whatever you draw last is
on top. Painting a picture, literally.

\begin{lstlisting}
DrawRectangle(640, 340, 100, 100, DARKPURPLE);   // painted first: underneath
DrawCircle(700, 400, 45, GOLD);                  // painted last:  on top
\end{lstlisting}

That means the structure of your draw section is your scene's depth. Background
first, then the world, then characters, then effects, then the HUD last of all.
When something is mysteriously invisible, nine times out of ten it is being drawn
\emph{before} something opaque.

\section{Text}

\begin{lstlisting}
DrawText("Hello", 20, 20, 20, DARKGRAY);        // text, x, y, fontSize, color
int width = MeasureText("Hello", 20);           // how wide would it be?
\end{lstlisting}

\lstinline|MeasureText| is how you centre text, since there is no alignment
parameter:

\begin{lstlisting}
const char *txt = "GAME OVER";
DrawText(txt, GetScreenWidth()/2 - MeasureText(txt, 60)/2, 200, 60, RED);
\end{lstlisting}

And to print variables, \lstinline|TextFormat| is raylib's
\lstinline|sprintf| that manages the buffer for you:

\begin{lstlisting}
DrawText(TextFormat("score: %i   time: %.2f", score, GetTime()), 20, 20, 20, BLACK);
\end{lstlisting}

\gotcha{\lstinline|TextFormat| writes into an internal rotating buffer. It is
perfect for drawing this frame and wrong for storing: never keep the returned
pointer across frames, copy the string if you need it later.}

\section{The shapes you actually use}

There are about forty drawing functions; these are the ones that cover almost
everything:

\begin{center}
\begin{tabular}{@{}ll@{}}
\toprule
\textbf{Function} & \textbf{Notes} \\
\midrule
\lstinline|DrawRectangle(x,y,w,h,c)| & corner based \\
\lstinline|DrawRectangleRec(rec,c)| & takes a \lstinline|Rectangle| \\
\lstinline|DrawRectangleLinesEx(rec,thick,c)| & outline with thickness \\
\lstinline|DrawRectangleRounded(rec,round,seg,c)| & rounded corners \\
\lstinline|DrawCircle(cx,cy,r,c)| & centre based \\
\lstinline|DrawCircleLines(cx,cy,r,c)| & outline \\
\lstinline|DrawLine(x1,y1,x2,y2,c)| & one pixel thin \\
\lstinline|DrawLineEx(v1,v2,thick,c)| & with thickness \\
\lstinline|DrawTriangle(v1,v2,v3,c)| & \textbf{counter-clockwise} vertices \\
\lstinline|DrawPoly(centre,sides,r,rot,c)| & regular polygon \\
\bottomrule
\end{tabular}
\end{center}

The naming pattern repeats across the whole library and is worth learning as a
rule rather than as forty separate facts:

\begin{itemize}[leftmargin=*,itemsep=2pt]
  \item plain name, e.g. \lstinline|DrawCircle| --- the simple version, loose
        \lstinline|int| coordinates;
  \item \textbf{V} suffix, e.g. \lstinline|DrawCircleV| --- same thing but taking
        a \lstinline|Vector2|;
  \item \textbf{Rec} suffix --- takes a \lstinline|Rectangle|;
  \item \textbf{Ex} suffix --- extended: extra parameters like thickness;
  \item \textbf{Pro} suffix --- all the knobs: origin, rotation, scaling;
  \item \textbf{Lines} in the name --- outline instead of filled.
\end{itemize}

Once you see that, you can guess most of raylib without looking anything up.

\gotcha{\lstinline|DrawTriangle| wants its three vertices in counter-clockwise
order. Give them clockwise and the triangle is facing away from you, so it gets
culled and you see nothing at all. If a triangle is invisible, swap two of its
vertices before you look for any other bug.}

\tryit{Lesson 2 draws the coordinate grid, every basic shape labelled, alpha
blending, and lets you swap the draw order with SPACE to watch the layering
change.}{./course2d/build.sh 2}
